\documentclass{cqutthesis}

\enablewatermark % 注释本行可编译无水印版

% 封面信息
\university{重庆理工大学}
% 题目一般不超过 20 个汉字，必要时可使用副标题。
\thesistitle{论文题目}
\college{学院全称}
\major{专业全称}
\classinfo{班级全称}
\authorname{学生姓名}
\studentid{学号}
\advisor{指导教师}
\advisortitle{职称}
\thesisdate{20XX 年 XX 月}

\begin{document}

\makecover
\makecommitment

% 目录和中英文摘要连续使用罗马数字页码，正文改用阿拉伯数字页码。
\frontmatterstart
\tableofcontents

\makeabstract{
% 中文摘要约 300 字，不放公式、图表和参考文献序号。
本文围绕论文研究对象，概述研究背景、研究方法、主要工作与研究结论。摘要应独立成文，内容简明、准确，能够反映论文的核心内容。
% 关键词填写 3--6 个，按外延由大到小排列，末尾不加标点。
}{关键词一；关键词二；关键词三}

\makeenglishabstract{
This abstract briefly presents the research background, methods, main work, and conclusions of the thesis. It should be self-contained and accurately reflect the core contribution of the study.
}{Keyword one, Keyword two, Keyword three}

\mainmatterstart

% 毕业论文正文原则上不少于 12000 字，各章均从新页开始。

\chapter{绪论}
\section{研究背景与意义}
在此撰写研究背景、研究目的及研究意义。正文采用统一的段落格式，文献引用使用上标方括号数字表示\cite{sample-journal}。

\section{国内外研究现状}
在此归纳国内外相关研究，并说明本文的研究切入点。

\section{研究内容与结构}
在此概述论文的研究内容、技术路线与章节安排。

\chapter{相关理论与研究方法}
\section{理论基础}
在此介绍与研究主题相关的理论基础。

\subsection{基本概念}
在此界定论文中的核心概念。

\subsection{研究方法}
在此说明所采用的研究方法、数据来源或实验设计。

\chapter{分析与讨论}
\section{研究结果}
在此呈现研究结果，并结合研究问题进行分析。

% 公式按章编号，例如本章首个公式为（3.1）。
\begin{equation}
  y = a + bx
\end{equation}

% 表题放在表格上方；学校规范采用三线表，表号形式为“表3-1”。
\begin{table}[htbp]
  \centering
  \caption{示例表题}
  \begin{tabular}{ccc}
    \toprule
    项目 & 指标一 & 指标二 \\
    \midrule
    示例 & 1 & 2 \\
    \bottomrule
  \end{tabular}
\end{table}

% 插图时将图题放在图片下方，图号形式为“图3.1”。
% \begin{figure}[htbp]
%   \centering
%   \includegraphics[width=0.7\textwidth]{assets/your-figure.png}
%   \caption{示例图题}
% \end{figure}

\section{讨论}
在此讨论研究结果的含义、适用范围与局限性。

\chapter{结论}
在此总结论文的主要结论，并提出后续研究展望。

\acknowledgements
感谢指导教师、参与研究的人员以及在论文完成过程中提供帮助的单位和个人。

\references
\begin{thebibliography}{99}
  % 毕业论文原则上不少于 15 篇，毕业设计不少于 10 篇；网络文献不计入上述数量。
  % 按正文首次引用顺序排列，著录格式执行 GB/T 7714—2015，每条末尾使用句号。
  \bibitem{sample-journal} 作者. 文献题名[J]. 期刊名称, 2026, 1(1): 1--10.
  \bibitem{sample-book} 作者. 著作名称[M]. 出版地: 出版社, 2026.
  \bibitem{sample-thesis} 作者. 学位论文题名[D]. 保存地: 保存单位, 2026.
\end{thebibliography}

\appendixchapter{补充材料}
在此放置不宜纳入正文、但对论文完整性有帮助的材料。理工科附录依次使用“附录1”“附录2”等编号。

\end{document}
